\chapter{Design and Evaluation of IPFS: A Storage Layer for the Decentralized Web}
\paperinfo{D. Trautwein, A. Raman, G. Tyson, I. Castro, W. Scott, M. Schubotz, B. Gipp, Y. Psaras (Protocol Labs, U. G\"ottingen, Telefonica Research, HKUST, QMUL, FIZ Karlsruhe)}{ACM SIGCOMM 2022}

\begin{keybox}[The paper in three sentences]
IPFS (InterPlanetary File System) is a peer-to-peer network for storing and fetching files in which a file is named by the
\key{hash of its content} (a \key{CID}), not by the server that holds it, so any peer can serve it and the receiver can check it.
To find who holds a CID, IPFS uses a \key{Kademlia DHT}: the holder publishes a \key{provider record} (CID $\to$ its PeerID) on the 20
peers whose IDs are closest to the CID, and a requester walks the DHT to find that record, then fetches the data with \key{Bitswap}. The paper
explains the design and measures the live network: about 200K peers in 152 countries and 2715 ASes, under 2.3\% in clouds, publication about 34\,s,
retrieval about 3\,s, and much faster through caching \key{HTTP gateways}.
\end{keybox}

% ------------------------------------------------------------------
\section{Motivation: the web has become centralised}

Today a new website typically runs on a cloud (AWS), uses a third-party DNS provider, a CDN (Akamai, Cloudflare) and a certificate authority
(Let's Encrypt). Each is well engineered, but each is a \key{single point of organisational failure}: when AWS, Cloudflare or OVH has an
outage, large parts of the web go down. The \key{Decentralized Web} movement wants web functions (naming, hosting, delivery) spread across
many independent operators. The most basic function every web application needs is \textbf{storing and serving objects}; IPFS aims to
decentralise exactly that, so other applications can be built on top.

\begin{keybox}[Four core ideas of IPFS]
\begin{enumerate}
  \item \key{Content addressing}: the name of an object is derived from its content, not its location, so it can be served by any peer.
  \item \key{Decentralised indexing}: a P2P DHT records who can provide each object; no central directory.
  \item \key{Immutability and self-certification}: the name \emph{is} a hash, so the receiver verifies data by re-hashing it; no certificates needed.
  \item \key{Open participation}: anyone can run a node, no permission needed.
\end{enumerate}
\end{keybox}

\textbf{HTTP vs IPFS in one line}: \texttt{https://example.com/cat.jpg} means ``ask \emph{that} server for whatever it calls cat.jpg''; \texttt{/ipfs/bafy\dots}
means ``give me the bytes whose hash is bafy\dots, from \emph{anyone}''.

% ------------------------------------------------------------------
\section{Building blocks}

\subsection{Content identifiers (CIDs)}

\paperfig[0.6]{ip_fig1}{Structure of a CIDv1: a multibase prefix, then the CID version, a multicodec (how the data is encoded), and a multihash (which hash function,
its length, and the hash itself).\label{fig:ipcid}}

A CID is \emph{self-describing}: it says how to read itself, so new hash functions or encodings can be added later without breaking old CIDs.
\begin{center}
\small
\begin{tabularx}{\linewidth}{@{}lX@{}}
\toprule
Field & Meaning (example in Figure~\ref{fig:ipcid})\\
\midrule
Multibase prefix & how the binary CID is written as text: \texttt{b} = base32 (24 encodings supported)\\
CID version & v1 (v0 also exists)\\
Multicodec & how the data is encoded: \texttt{dag-pb} (protobuf), json, cbor, \dots\\
Multihash & hash function (\texttt{sha2-256}), digest length (32 bytes), then the digest itself\\
\bottomrule
\end{tabularx}
\end{center}

\textbf{How a file becomes a CID}:
\begin{enumerate}
  \item Split the file into \key{chunks} (default 256\,kB). Each chunk's CID is the hash of its bytes (plus the metadata above).
  \item Build a \key{Merkle DAG}: a parent node lists its children's CIDs, and its own CID is the hash of that list. The top node's CID is the \key{root CID}, the name of
        the whole file.
\end{enumerate}

\textbf{Example.} A 600\,kB file becomes chunks A (256\,kB), B (256\,kB), C (88\,kB) with CIDs $c_A, c_B, c_C$. The root node contains $[c_A, c_B, c_C]$, and the file's
name is root $= H([c_A, c_B, c_C])$.
\begin{itemize}
  \item \textbf{Verification}: when you receive chunk B, hash it and compare with $c_B$. If any byte was changed, the hash differs. So you can fetch chunks from
        untrusted strangers.
  \item \textbf{De-duplication}: a Merkle DAG (unlike a tree) lets a node have several parents. If a second file contains the same chunk B, it just points to $c_B$; B is
        stored and sent once.
  \item \textbf{Location independence}: the DAG says nothing about where chunks live, so copying or deleting replicas never changes any CID.
  \item \textbf{The catch: immutability}. Edit one byte and the root CID changes. For content that changes (a blog's front page) IPFS has \key{IPNS}: a name derived from
        the publisher's public key, which points to the current CID via a record \emph{signed} with the private key. The publisher updates the record; the name stays the same.
\end{itemize}

\subsection{Peer identities and Multiaddresses}
\begin{itemize}
  \item Each peer generates a public/private key pair. Its \key{PeerID} is the hash of its public key. When two peers open a secure channel, the PeerID proves the other side
        owns that key.
  \item A peer's location is a \key{Multiaddress}, a self-describing path of protocol layers, e.g.\ \texttt{/ip4/1.2.3.4/tcp/3333/p2p/QmZyWQ14\dots}: IPv4 address,
        TCP port, then the PeerID. Why this format: (1) different peers speak different transports (TCP, QUIC, WebSocket), and the address says up front whether you can
        reach it; (2) addresses can be \emph{stacked} to say ``reach X through relay Y'', which is how browser nodes (which cannot accept connections) are reached.
\end{itemize}

\subsection{The DHT: Kademlia, as used by IPFS}
To find who has a CID, IPFS stores mappings in a \key{Kademlia} DHT (the XOR geometry from the DHT-geometry chapter). CIDs and PeerIDs are both hashed with SHA-256 into
one \textbf{256-bit key space}.
\begin{itemize}
  \item \textbf{Distance} between two keys = their bitwise \key{XOR}, read as a number. Smaller XOR = ``closer''.
  \item \textbf{Routing table}: 256 \key{$k$-buckets}; bucket $i$ holds up to $k = 20$ peers whose distance from me is in $[2^i, 2^{i+1})$. So I know many peers near me and a few
        in each far region.
  \item \textbf{Lookup}: iterative. Ask the $\alpha = 3$ closest peers I know; each replies with either the record or closer peers it knows; repeat with the closest new ones.
\end{itemize}
\textbf{XOR example (4-bit IDs).} Key $K = 1011$. Peers: $P_1 = 1001$, $P_2 = 0011$, $P_3 = 1010$.
\begin{center}
\begin{tabular}{@{}lll@{}}
\toprule
Peer & $P \oplus K$ & Distance\\
\midrule
$P_1 = 1001$ & 0010 & 2\\
$P_2 = 0011$ & 1000 & 8\\
$P_3 = 1010$ & 0001 & \textbf{1 (closest)}\\
\bottomrule
\end{tabular}
\end{center}
Peers that share a long prefix with the key are close. A record for $K$ is stored on the peers with the smallest XOR distance.

\textbf{IPFS's changes to textbook Kademlia}: SHA-256 (256-bit) keys instead of SHA-1 (160-bit), to stay safe as SHA-1 collisions become feasible; $k = 20$ per bucket;
and TCP/QUIC instead of UDP, which made connection management simpler.

\subsection{DHT servers and DHT clients (Autonat)}
Many peers sit behind NATs and cannot accept incoming connections. If such a peer were in others' routing tables, every lookup through it would wait for a timeout. So:
\begin{enumerate}
  \item A new peer joins as a \key{DHT client}.
  \item It asks other peers to \emph{dial it back} (\key{Autonat}).
  \item If \textbf{more than 3} succeed, it is publicly reachable and becomes a \key{DHT server}: it stores and serves records. Otherwise it stays a client: it can make
        requests, but does not store records and is kept out of routing tables.
\end{enumerate}
The authors credit this split (from v0.5) with a large speed-up, since lookups no longer waste time on unreachable peers.

% ------------------------------------------------------------------
\section{Publishing and retrieving content}

\paperfig[0.55]{ip_fig3}{Publication (steps 1--3) and retrieval (steps 4--6) in IPFS.\label{fig:ipflow}}

\subsection{Publication (steps 1--3)}
\begin{enumerate}
  \item[(1)]\textbf{Import}: the provider adds the file to its \emph{local} IPFS node; it is chunked and gets a root CID. Nothing is uploaded anywhere: the data
        stays on the provider.
  \item[(2)]\textbf{DHT walk}: find the \key{20 peers whose PeerIDs are closest (by XOR) to the CID's hash}.
  \item[(3)]\textbf{Store the provider record} (CID $\to$ my PeerID) on those 20 peers, by RPC, ``fire and forget'' (does not wait for every reply).
\end{enumerate}
Separately, each peer publishes a \key{peer record} (PeerID $\to$ its Multiaddresses) in the same way, so others can turn a PeerID into an IP address.

\textbf{Record lifetimes}:
\begin{itemize}
  \item \textbf{Replication $k = 20$}: peers come and go (churn), so a record on one peer would soon vanish. 20 is a compromise between overhead and the risk of losing the record.
  \item \textbf{Republish every 12\,h}: if the original 20 holders have left, the provider picks the current 20 closest.
  \item \textbf{Expire after 24\,h}: if the provider stops republishing (it went offline), its record disappears, so nobody is sent to a dead provider.
\end{itemize}
Peers behind NATs cannot serve content, so they pay \key{pinning services} to host and publish it for them. Any peer that downloads content may also announce itself as a
provider; downloaders do not need to trust it, since they verify the hash.

\subsection{Retrieval (steps 4--6)}
Four logical phases:
\begin{enumerate}
  \item[(4)]\textbf{Opportunistic Bitswap}: first ask all peers I am \emph{already connected to} whether they have the CID. Neighbours often do, and this avoids the DHT.
        If nobody answers within \key{1\,s}, fall back to the DHT.
  \item[(5)]\textbf{Content discovery}: a DHT walk to find a provider record for the CID (stops at the first record found).
  \item[(5b)] \textbf{Peer discovery}: a \emph{second} DHT walk to turn the provider's PeerID into a Multiaddress (skipped if it is in my address book of about 900 recently seen peers).
  \item[(6)]\textbf{Peer routing and content exchange}: connect, and fetch with \key{Bitswap}: send \texttt{WANT-HAVE} (``do you have this block?''), the peer answers
        \texttt{HAVE}, then send \texttt{WANT-BLOCK} and receive the block. Verify each block against its CID.
\end{enumerate}

\textbf{Same CID, two requesters}:
\begin{center}
\small
\begin{tabularx}{\linewidth}{@{}lXX@{}}
\toprule
Step & Alice (her neighbour already has it) & Bob (no neighbour has it)\\
\midrule
Bitswap to neighbours & neighbour says HAVE: fetch, verify, done & silence for 1\,s\\
DHT walk 1 (provider) & -- & finds record CID $\to$ PeerID$_P$\\
DHT walk 2 (address) & -- & finds PeerID$_P$ $\to$ /ip4/\dots/tcp/\dots\\
Connect + Bitswap & -- & fetch, verify\\
Afterwards & may announce herself as a provider & may announce himself as a provider\\
\bottomrule
\end{tabularx}
\end{center}

\subsection{HTTP gateways}
Most users do not run IPFS. A \key{gateway} is a bridge: one side is a DHT-server IPFS node, the other an nginx web server that accepts
\texttt{GET https://gateway/ipfs/\{CID\}}. Two caches make it fast:
\begin{itemize}
  \item the \key{nginx cache} (LRU) of recently served objects;
  \item the gateway's own \key{IPFS node store}, holding content \emph{pinned} there (e.g.\ by the web3.storage and nft.storage initiatives).
\end{itemize}
Gateways are optional and anyone can run one (107 were known). They are a centralising element, but not a required one.

% ------------------------------------------------------------------
\section{How the network was measured}
No one sees all of IPFS, so the authors combine three datasets:
\begin{enumerate}
  \item \textbf{Crawler} (active): every 30 minutes, start from the 6 bootstrap peers and recursively ask every peer for all entries of its $k$-buckets until no new peers appear.
        Map IPs to country and AS. Revisit peers to measure session lengths (probe interval half the observed uptime, between 30\,s and 15\,min).
  \item \textbf{Gateway logs} (passive): one day of the \texttt{ipfs.io} gateway, 7.1M requests.
  \item \textbf{Controlled benchmarks}: six AWS VMs (Cape Town, Sydney, Frankfurt, Bahrain, S\~ao Paulo, N.\ California) running IPFS DHT servers in the \emph{public} network.
        Each round, one publishes a new 0.5\,MB object and the other five retrieve it, then disconnect so the next round cannot use Bitswap and must use the DHT.
\end{enumerate}

% ------------------------------------------------------------------
\section{What the network looks like}
\begin{itemize}
  \item \textbf{Size}: about \key{199K PeerIDs}, 464K IP addresses, \key{152 countries}; 45.5\% of IPs were never reachable. Top countries: US 28.5\%, China 24.2\%.
  \item \textbf{Reliable peers} ($>90\%$ uptime) are only 1.4\%, spread evenly by country: no single country could take IPFS down.
  \item \textbf{Peers per IP}: 92.3\% of IPs host one PeerID, but the top 10 IPs host about 66K PeerIDs between them, a possible risk (a misbehaving host could rotate IDs).
  \item \textbf{ASes}: \key{2715 ASes}, but concentrated: the top 10 ASes hold 64.9\% of IPs, and two Chinese ASes alone over 30\%.
  \item \textbf{Cloud}: under \key{2.3\%} of nodes run in cloud providers (compare Mastodon, with 6\% on Amazon alone): most operators run their own hardware.
  \item \textbf{Churn}: 87.6\% of sessions last under 8\,h; only 2.5\% exceed 24\,h; median uptime in Hong Kong is 24\,min. This is why records are replicated on 20 peers and
        republished every 12\,h.
\end{itemize}

% ------------------------------------------------------------------
\section{Performance}

\subsection{Publication: slow, because of the DHT walk}
Median \key{33.8\,s}; 90th percentile 112\,s; 95th 138\,s; similar in every region, and independent of file size (only the record is published).
\begin{itemize}
  \item The \key{DHT walk} to find the 20 closest peers takes \key{87.9\%} of the time.
  \item The RPCs to store records show spikes at protocol timeouts: 5\,s (TCP/QUIC dial timeout) and 45\,s (WebSocket handshake timeout), caused by unresponsive peers.
\end{itemize}

\subsection{Retrieval: seconds}
100\% of retrievals succeeded. Median \key{2.90\,s}, 90th 4.34\,s, 95th 4.74\,s. By region the median ranges from \key{1.81\,s} (Frankfurt) to 3.75\,s (Cape Town).
\begin{itemize}
  \item Every retrieval in this setup includes the \key{1\,s Bitswap timeout} (the experiment ensured neighbours never had the data): a worst case.
  \item A retrieval DHT walk has median \key{622\,ms}, much faster than a publication walk, because it stops at the first record instead of finding 20 peers. Both walks
        together are under 2\,s for half the retrievals. (Earlier Kademlia deployments in BitTorrent took over a minute per lookup, mainly because of dead nodes.)
  \item Fetching the 0.5\,MB itself takes under 1.26\,s in 99\% of cases.
\end{itemize}

\subsection{The cost of decentralisation: stretch}
\[
\text{Stretch} = \frac{\text{IPFS retrieval time}}{\text{HTTPS retrieval time}} = \frac{\text{Discover} + \text{Dial} + \text{Negotiate} + \text{Fetch}}{\text{Dial} + \text{Negotiate} + \text{Fetch}}
\]
Dial $\approx$ TCP handshake, Negotiate $\approx$ TLS handshake, Fetch $\approx$ GET and transfer. HTTPS needs no Discover step, because the URL already says where the content is.

\textbf{Example.} Discover $= 1.0$\,s (Bitswap timeout) $+ 1.2$\,s (two DHT walks) $= 2.2$\,s; Dial 0.1\,s; Negotiate 0.1\,s; Fetch 0.4\,s.
\begin{center}
\begin{tabular}{@{}lcc@{}}
\toprule
 & With Bitswap timeout & Without it\\
\midrule
IPFS time & $2.2 + 0.6 = 2.8$\,s & $1.2 + 0.6 = 1.8$\,s\\
HTTPS time & 0.6\,s & 0.6\,s\\
Stretch & \textbf{4.7} & \textbf{3.0}\\
\bottomrule
\end{tabular}
\end{center}
Measured: median stretch \key{4.3}. Without the 1\,s Bitswap wait, Frankfurt gets stretch $< 2$ for 80\% of retrievals. Lesson: \key{run the DHT lookup in parallel with
Bitswap} instead of after it; this costs extra messages but saves the 1\,s.

\subsection{Gateways: caching makes it fast}
\begin{center}
\begin{tabular}{@{}lccc@{}}
\toprule
Served from & Median latency & Share of requests & Share of bytes\\
\midrule
nginx cache & 0\,s & 46.0\% & 46.4\%\\
gateway's IPFS node store (pinned) & 8\,ms & 40.2\% & 38.0\%\\
not cached (fetched from IPFS) & 4.04\,s & 13.8\% & 15.6\%\\
\bottomrule
\end{tabular}
\end{center}
\begin{itemize}
  \item Combined hit rate is over \key{80\%}, and \key{76\%} of requests are served in under 250\,ms.
  \item Object size does not predict latency (correlation 0.13): delay comes from size-independent steps like DHT walks. Median object 665\,KB.
  \item 101K users in one day; 51.8\% of traffic was referred by ordinary websites, many of them video and NFT sites, which store NFT data on IPFS.
\end{itemize}

% ------------------------------------------------------------------
\section{Summary comparison}
\begin{center}
\small
\begin{tabularx}{\linewidth}{@{}lXX@{}}
\toprule
 & HTTP/Web today & \textbf{IPFS}\\
\midrule
Name means & a location (server + path) & the content's hash (CID)\\
Who can serve & that server (or its CDN) & \key{any} peer with the bytes\\
Integrity & trust the server via TLS certificates & \key{self-certifying}: re-hash and compare\\
Finding the content & DNS lookup & Bitswap to neighbours, then Kademlia DHT walk(s)\\
Mutable content & natural & needs IPNS (signed pointer)\\
Typical latency & sub-second & seconds (median stretch 4.3); gateway cache $<$250\,ms for 76\%\\
\bottomrule
\end{tabularx}
\end{center}

\begin{keybox}[Exam one-liner]
IPFS is a content-addressed P2P storage network: files are chunked (256\,kB) into a Merkle DAG whose root hash is the self-describing CID (multibase, version, multicodec,
multihash), giving verification, de-duplication and location independence (IPNS handles mutable content). Peers are named by PeerID (hash of public key) and located by
Multiaddresses. A Kademlia DHT (256-bit XOR space, $k=20$, $\alpha=3$; only publicly reachable peers become DHT servers via Autonat) stores provider records on the 20 closest
peers (republish 12\,h, expire 24\,h) and peer records. Retrieval: Bitswap to neighbours (1\,s timeout) $\to$ DHT walk for the provider $\to$ DHT walk for its address $\to$ Bitswap
fetch. Measured: about 200K peers, 152 countries, 2715 ASes, $<$2.3\% cloud, heavy churn; publication about 34\,s (DHT walk 88\%), retrieval about 3\,s (stretch 4.3), gateways with
caching serve 76\% under 250\,ms.
\end{keybox}

\section{Revision checklist}
\begin{itemize}
  \item Why: centralisation $\to$ single points of organisational failure; four core ideas.
  \item CID fields; chunking; Merkle DAG; verification, de-duplication, location independence; immutability and IPNS.
  \item PeerID = hash(public key); Multiaddress layers and relaying.
  \item Kademlia: XOR distance (do an example), $k$-buckets, iterative lookup with $\alpha = 3$; IPFS changes (SHA-256, $k=20$, TCP/QUIC).
  \item DHT server vs client via Autonat (more than 3 successful dial-backs).
  \item Publication: import, DHT walk to 20 closest, store provider record; peer record; republish 12\,h, expire 24\,h; pinning services for NATed users.
  \item Retrieval: Bitswap first (1\,s), DHT walk for provider, DHT walk for address (address book of 900), Bitswap WANT-HAVE/HAVE/WANT-BLOCK.
  \item Gateways: nginx LRU cache + pinned node store.
  \item Measurement: crawler, gateway logs, 6 AWS vantage points.
  \item Numbers: 199K peers, 152 countries, 2715 ASes (top 10 = 65\%), $<$2.3\% cloud, 2.5\% of sessions $>$24\,h; publication 33.8\,s median; retrieval 2.9\,s median; DHT walk 622\,ms; stretch 4.3; gateway hit rate $>$80\%.
\end{itemize}
